\documentclass{article}
\usepackage[T1]{fontenc}
\usepackage{lmodern}
\usepackage{graphicx}
\usepackage{amsmath,amssymb}
\usepackage[a4paper,margin=2cm]{geometry}
\usepackage[absolute,overlay]{textpos}
\setlength{\TPHorizModule}{1mm}
\setlength{\TPVertModule}{1mm}
\usepackage[indonesian]{babel}
\usepackage{hyperref}
\setlength{\emergencystretch}{2em}
% unit: Halaman 8

\begin{document}

% === Halaman 8 (llm) ===

\begin{itemize}
    \item Solusi dari \textit{superincreasing knapsack} (yaitu $b_1, b_2, ..., b_n$) mudah dicari sebagai berikut (berarti sama dengan mendekripsikan cipherteks menjadi plainteks semula):
\end{itemize}

\begin{enumerate}
    \item Jumlahkan semua bobot di dalam barisan.
    \item Bandingkan bobot total dengan bobot terbesar di dalam barisan. Jika bobot terbesar lebih kecil atau sama dengan bobot total, maka ia dimasukkan ke dalam \textit{knapsack}, jika tidak, maka ia tidak dimasukkan.
    \item Kurangi bobot total dengan bobot yang telah dimasukkan, kemudian bandingkan bobot total sekarang dengan bobot terbesar selanjutnya. Demikian seterusnya sampai seluruh bobot di dalam barisan selesai dibandingkan.
    \item Jika bobot total menjadi nol, maka terdapat solusi persoalan \textit{superincreasing knapsack}, tetapi jika tidak nol, maka tidak ada solusinya.
\end{enumerate}

\end{document}
